\textbf{Задача 21} Докажите, что имея источник случайных битов, распределенных независимо и равных единице с вероятностью $p$, можно смоделировать случайные биты, распределённые равномерно и независимо. Какое среднее число неравномерных битов нужно для моделирования одного равномерного?

\begin{proof} Возьмём два неравномерных бита. Если они не равны друг другу, то возьмём первый из них в качестве равномерного. Если же равны, возьмём ещё два неравномерных бита, и так далее. Вероятность того, что выпали разные биты, равна $2p(1 - p)$, причём условные вероятности комбинаций 01 и 10 равны по $\frac12$. \\
\\
Вероятность того, что выпадут два одинаковых неравномерных бита равна $1-2p+2p^2$. Тогда среднее число неравномерных битов, нужных для получения одного случайного равномерного, равно $2\mathbb{E}[$число шагов до получения двух различных неравномерных битов$] = 2\Sigma_{k=1}^\infty k (1-2p+2p^2)^{k-1}(2p-2p^2) = \frac{1}{p(1-p)}$  \end{proof}